\documentclass[11pt]{article}
\usepackage[T1]{fontenc}
\usepackage{lmodern,amsmath,amssymb,amsthm,booktabs}
\usepackage[letterpaper,margin=1in]{geometry}
\usepackage{microtype,xcolor}
\usepackage[colorlinks=true,linkcolor=blue!45!black,urlcolor=blue!45!black]{hyperref}
\newtheorem{theorem}{Theorem}[section]
\newtheorem{lemma}[theorem]{Lemma}
\newtheorem{proposition}[theorem]{Proposition}
\newtheorem{corollary}[theorem]{Corollary}
\newtheorem{remark}[theorem]{Remark}
\newcommand{\PP}{\mathbb P}
\newcommand{\EE}{\mathbb E}
\newcommand{\Z}{\mathcal Z}
\newcommand{\cO}{\mathcal O}
\newcommand{\comp}{\operatorname{comp}}
\newcommand{\tr}{\operatorname{tr}}
\newcommand{\half}{\tfrac12}
\setlength{\emergencystretch}{3em}
\title{A stronger unrestricted second-order cycle bound}
\author{Research working draft}
\date{10 November 2026}
\begin{document}
\maketitle
\begin{abstract}
Let $\psi(r)$ be the maximum number of cycles in a finite graph of cycle rank $r$.
We prove that, for an absolute constant $M$,
\[
 \psi(r)\le2^{r-1}+(M+o(1))\frac{2^r(\log_2 r)^{11/2}}{\sqrt r}.
\]
In particular, every fixed exponent $\kappa<1/2$ is valid in a relative error $O(r^{-\kappa})$.
The proof combines an exponential cut estimate, independence of parallel link components, and two nested short-cycle deletions.
A weighted separation argument produces a long chain of nested cycles.
Restoring the medium-length deleted cycles leaves many regions with small boundary.
Each region supplies a positive amount of cycle weight, contradicting a second-moment upper bound.
The proposed error $O(r2^{r/2})$ remains open without structural restrictions.
\end{abstract}

\section*{Status and contributions}
This draft extends the existing cycle-space machinery in the shared research programme.
Codex derived the exponential cut estimates and the endpoint parameter refinement.
Fable 5.2 supplied the nested-cycle argument, the independent parallel-link lemma, and the simplification to a separate spectral estimate in each restored block.
The final proof combines these contributions.
The earlier two-scale proof is retained to make the geometry and the direct cycle transfer explicit.
The new endpoint proof and its sharper packing refinement have passed independent mathematical audits.
Publication priority has not been established.

Fable formalized the parallel-link lemma, the link bound, separator geometry, averaging, the refined moment, and the transfer from a forest estimate in Lean.
An independent build of these modules succeeded.
The global forest estimate and the assembly of the full theorem are not yet formalized.
No existing root Lean import or original manuscript was changed by this research pass.

\section*{Conventions and reading guide}
Rank means $|E|-|V|+\operatorname{comp}$.
An even subgraph is an edge set with even degree at every vertex.
Graphs may have loops and parallel edges; distinct edges remain distinct.
A loop is a cycle of length one, and a parallel pair is a cycle of length two.
Subdivision transfers cycle counts and rank to simple graphs.

The first three sections prove the exponential cut estimate used in the main argument.
The next section supplies the cycle-space, walk, and spectral identities.
The remaining sections give the first polynomial proof and its cycle transfer, then prove the stronger endpoint estimate.
The last theorem uses the stronger branch of the cut estimate to obtain the logarithmic power $11/2$.
The direct cubic-extremizer transfer gives a self-contained route to unrestricted cycle counts.
The intermediate cut bounds are retained to make the proof self-contained.

\clearpage
\tableofcontents
% New research fragment. Root and path_stability independently audited the safe exponential bound.
% Do not merge into the previously audited manuscripts without a separate review.
\section{Exponential decay for disjoint regions with bounded cuts}
\label{sec:second-order-exponential-cuts}

All cycle spaces below are over $\mathbb F_2$.
Graphs may have parallel edges.
For a graph $Q$, write $\mathcal Z(Q)$ for its binary cycle space.
A uniform element of this space is denoted by $F$.
A vertex is \emph{closed} when no incident edge belongs to $F$.
Loops can be deleted when discussing this event, since it concerns the edge boundary of a vertex.

\begin{lemma}[A common local functional]\label{lem:common-star-functional}
Let $Q$ be a connected loopless multigraph.
Let $v_1,\ldots,v_\ell$ be distinct vertices of degree at most $d$, where $d\ge2$.
Suppose a nonzero functional $a\in\mathcal Z(Q)^*$ lies in the span of the incident edge-coordinate functionals at every $v_i$.
If $Z$ counts the closed vertices among $v_1,\ldots,v_\ell$, then
\[
 \mathbb E\,2^{-Z}
 \le\frac12+\frac12\left(1-2^{1-d}\right)^{\lfloor\ell/2\rfloor}.
\]
\end{lemma}

\begin{proof}
At $v_i$, choose an incident-edge subset $A_i$ whose coordinate sum represents $a$.
It is nonempty and proper: the empty subset and the whole star both represent zero.
Split $v_i$ into two vertices carrying $A_i$ and its complement, and join them by a new edge $h_i$.
Perform all splits simultaneously, obtaining a connected graph $Q'$.
Contracting the new edges identifies the two cycle spaces.
Under this identification, the coordinate on every $h_i$ equals $a$.

Deleting all $\ell$ new edges imposes exactly the one condition $a=0$.
Thus
\[
 \beta\bigl(Q'-\{h_1,\ldots,h_\ell\}\bigr)=\beta(Q')-1.
\]
Keeping isolated vertices in the rank calculation shows that the remaining graph has exactly $\ell$ components.
Contract these components.
The resulting connected quotient has $\ell$ vertices and $\ell$ edges.
An even subgraph with $a=1$ projects to its entire edge set.
The quotient is therefore connected and two-regular, hence a cycle.
For $\ell=1$ it is a loop, and for $\ell=2$ it consists of two parallel edges.

Condition on $a=0$.
The restrictions to the $\ell$ components are independent uniform even subgraphs.
Choose a matching of $\lfloor\ell/2\rfloor$ edges in the quotient cycle.
For each chosen $h_i$, its two split endpoints belong to distinct components.
Different chosen edges use disjoint component pairs.

Suppose the two parts of the original star have sizes $s_i$ and $t_i$.
Their sum is $\deg_Q(v_i)\le d$.
In its component, a vertex of degree $s_i$ has at most $s_i-1$ independent incident-edge coordinates.
Its star is therefore zero with probability at least $2^{1-s_i}$.
The corresponding bound at the other endpoint is $2^{1-t_i}$.
Consequently, the original vertex $v_i$ is closed with conditional probability at least $2^{2-d}$.
For the chosen matching, these events are independent.
It follows that
\[
 \mathbb E[2^{-Z}\mid a=0]
 \le\left(1-2^{1-d}\right)^{\lfloor\ell/2\rfloor}.
\]
When $a=1$, none of the selected original vertices is closed.
Since $a$ is a nonzero functional, its two values have equal probability.
This proves the result.
For $\ell=1$, the displayed bound is the trivial bound one.
\end{proof}

\begin{theorem}[Exponential closed-star estimate]\label{thm:exponential-closed-stars}
Let $Q$ be a connected loopless multigraph with $k\ge1$ distinct marked vertices of degree at most $d$.
Let $Z$ count the marked vertices closed in a uniform even subgraph.
For $d\ge2$, put
\[
 T_d=2^d,\qquad D_d=(d-1)(2^d-1),\qquad M_d=2^{D_d}-1.
\]
Then
\begin{equation}\label{eq:exponential-closed-stars}
 \mathbb E\,2^{-Z}
 \le\frac12+\exp\!\left(-\frac{k}{2^d M_d}\right)
 \le\frac12+\exp\!\left(-\frac{k}{2^{D_d+d}}\right).
\end{equation}
For $d\le1$, every marked vertex is closed and $\mathbb E2^{-Z}=2^{-k}$.
\end{theorem}

\begin{proof}
Let $L_i\subseteq\mathcal Z(Q)^*$ be the span of the edge-coordinate functionals at marked vertex $i$.
The parity equation at that vertex gives $\dim L_i\le d-1$.
If some $L_i$ is zero, that vertex is always closed.
The expectation is then at most one half, so assume every $L_i$ is nonzero.

If $T_d$ of these spaces form a direct sum, their local patterns are independent.
Each corresponding vertex is closed with probability at least $2^{1-d}$.
Consequently,
\[
 \mathbb E2^{-Z}\le(1-2^{-d})^{T_d}\le e^{-1}<\frac12.
\]

Otherwise choose a maximal direct-sum family and let $U$ be its span.
It has fewer than $T_d$ members, so $\dim U\le D_d$.
Every unselected $L_i$ meets $U$ nontrivially, by maximality.
Every selected $L_i$ also meets $U$ nontrivially, since it is nonzero and contained in $U$.
There are at most $M_d$ nonzero vectors in $U$.
Thus one nonzero functional $a$ belongs to at least $\ell\ge k/M_d$ of the spaces $L_i$.

Apply Lemma~\ref{lem:common-star-functional} to those vertices and ignore the other marked factors.
With $q=1-2^{1-d}\ge1/2$, one has
\[
 \frac12q^{\lfloor\ell/2\rfloor}
 \le\frac1{2q}\exp\!\left(-\frac{(1-q)\ell}{2}\right)
 \le\exp(-2^{-d}\ell).
\]
This proves~\eqref{eq:exponential-closed-stars}.
The assertion for $d\le1$ follows directly from parity.
\end{proof}

\begin{corollary}[Disjoint cyclic regions]\label{cor:exponential-small-cuts}
Let $G$ be a connected graph and let $U_1,\ldots,U_k$ be pairwise disjoint vertex sets.
Suppose each $G[U_i]$ is connected and contains a cycle, and $|\delta_G(U_i)|\le d$.
For a uniform even subgraph $F$ of $G$, let $\mathcal F$ be the event that every $F\cap E(G[U_i])$ is a forest.
For $d\ge2$,
\[
 \Pr(\mathcal F)
 \le\frac12+\exp\!\left(-\frac{k}{2^d(2^{(d-1)(2^d-1)}-1)}\right).
\]
For $d\le1$, one has $\Pr(\mathcal F)\le2^{-k}$.
\end{corollary}

\begin{proof}
Contract each connected induced region to a marked vertex and delete the resulting internal loops.
Retain all other vertices and edges, obtaining a connected quotient $Q$.
Projection maps $\mathcal Z(G)$ onto $\mathcal Z(Q)$.
Its kernel is the direct sum of the cycle spaces of the induced regions.
Thus the quotient pattern is uniform.
Conditional on that pattern, the internal choices in different regions are independent affine cycle-space choices.

If a marked quotient vertex is closed, the corresponding internal choice is a uniform even subgraph of $G[U_i]$.
Its only forest is the empty edge set.
Since that region has positive rank, its forest probability is at most one half.
For a nonclosed marked vertex, bound the conditional forest probability by one.
If $Z$ is the number of closed marked vertices, these observations give
\[
 \Pr(\mathcal F\mid\text{quotient pattern})\le2^{-Z}.
\]
Any remaining quotient loops are at unmarked vertices.
Delete them before applying the loopless closed-star theorem.
Their independent coordinates do not change any marked closure event.
Average and apply Theorem~\ref{thm:exponential-closed-stars}.
\end{proof}

\paragraph{Quantitative scope.}
For each fixed cut bound $d$, the error decreases exponentially with the number of disjoint cyclic regions.
The explicit constant is large: its binary logarithm is $O(d2^d)$.
This theorem improves the previous inverse-linear dependence on the region count.
It does not by itself establish a polynomial error for unrestricted cycle counts.
When applying it with a growing cut bound, the dependence on $d$ must be retained.

% Independent research fragment. Root derivation; independent audits passed.
\section{A terminal-forest improvement to the cut constant}
\label{sec:terminal-forest-cut-constant}

The common-functional lemma gives a much better constant when the local spaces come from a graph.
The following argument replaces a vector-space pigeonhole by a forest with few terminals.

\begin{lemma}[A local intersection in a terminal forest]
\label{lem:local-terminal-intersection}
Let $Q$ be a connected loopless multigraph, and let $S$ be a nonempty proper vertex set.
Write $L_v$ for the span of the edge-coordinate functionals at $v$ in $\mathcal Z(Q)^*$.
Put $U=\sum_{v\in S}L_v$ and $B=|\delta_Q(S)|$.
Let $\mathcal V$ be a set of vertices outside $S$ with $L_v\cap U\ne0$.
Except for at most $2B$ members, $\mathcal V$ can be partitioned into at most $2B$ classes.
Each class has a common nonzero functional in all its spaces $L_v$.
\end{lemma}

\begin{proof}
At most $B$ vertices outside $S$ have a neighbor in $S$.
Exclude those vertices.
Choose a spanning tree in each component of $Q-S$.
For each edge of $\delta_Q(S)$, attach a distinct terminal leaf to its endpoint outside $S$.
Let $T$ be the union of paths between terminal leaves in these trees.
Components with fewer than two terminals contribute nothing.
The forest $T$ has at most $B$ leaves.
It has at most $B$ vertices of degree at least three.
After suppressing its degree-two vertices, it has at most $2B$ edges.
Exclude the members of $\mathcal V$ at branching vertices of $T$.

We show that every remaining member lies internally on one of these edges.
Fix $v\in\mathcal V$ with no neighbor in $S$, and choose $0\ne a\in L_v\cap U$.
Represent $a$ both by a subset $A$ of the star at $v$ and by a subset of edges incident with $S$.
The symmetric difference of these two representing sets annihilates $\mathcal Z(Q)$.
It is therefore a cut of $Q$.
Let $H$ be the component of $Q-S$ containing $v$.
Every component of $H-v$ lies entirely on one side of this cut.
Consequently, $A$ consists of the edges from $v$ to a union of components of $H-v$.
As a functional, $a$ equals the sum of the boundary-terminal coordinates attached to that union.

The sum of all terminal coordinates of $H$ is zero, because $\delta_Q(H)$ is a cut.
Since $a\ne0$, at least two components of $H-v$ contain terminals.
Every spanning tree of $H$ must connect these terminal groups through $v$.
Thus $v$ has degree at least two in $T$.
It has degree exactly two after the branching exceptions have been removed.
There are exactly two terminal-bearing components of $H-v$.
Their terminal sets are the two sides of the corresponding edge of the suppressed forest.
The sum of the terminal coordinates on either side is $a$.
In particular, this common functional is nonzero.

All remaining members on the same suppressed edge give the same terminal bipartition.
They therefore contain the same nonzero functional in their local spaces.
The adjacent and branching exceptions number at most $2B$.
The number of classes is at most $2B$, as required.
\end{proof}

\begin{theorem}[An explicit exponential cut estimate]
\label{thm:terminal-forest-cut-estimate}
Let $Q$ be a connected loopless multigraph with $k$ distinct marked vertices of degree at most $d\ge2$.
Let $Z$ count the marked vertices closed in a uniform even subgraph.
Then
\[
 \mathbb E2^{-Z}\le\frac12+2\exp\!\left(-\frac{k}{2d4^d}\right).
\]
The same bound applies to the probability that a uniform even subgraph is a forest in each of $k$ disjoint cyclic regions.
Each region must have an edge boundary of size at most $d$.
\end{theorem}

\begin{proof}
Use the local spaces $L_v$ from Theorem~\ref{thm:exponential-closed-stars}.
A zero local space gives an expectation at most one half.
If $2^d$ local spaces form a direct sum, that theorem's independence argument also gives an expectation below one half.

Otherwise take a maximal direct-sum family, with vertex set $S$ and size $t<2^d$.
Every remaining local space meets $U=\sum_{v\in S}L_v$ nontrivially.
Suppose first that $k>(1+2d)2^d$.
Then $S$ is nonempty and proper, and $B=|\delta_Q(S)|\le dt$.
Lemma~\ref{lem:local-terminal-intersection} gives a common nonzero functional for at least
\[
 \ell\ge\frac{k-t-2B}{2B}
 \ge\frac{k}{2d2^d}-2
\]
marked vertices.
The common-functional lemma, with $q=1-2^{1-d}\ge1/2$, gives
\[
 \mathbb E2^{-Z}-\frac12
 \le\frac12q^{\lfloor\ell/2\rfloor}
 \le e^{-2^{-d}\ell}
 \le e^{2^{1-d}}\exp\!\left(-\frac{k}{2d4^d}\right)
 \le2\exp\!\left(-\frac{k}{2d4^d}\right).
\]
If $k\le(1+2d)2^d$, the asserted right side exceeds one, so the result is trivial.
The quotient argument in Corollary~\ref{cor:exponential-small-cuts} transfers this estimate to cyclic regions.
\end{proof}

\paragraph{Status and scope.}
The proof has passed independent mathematical audits, including an audit by Fable 5.2.
It is not formalized in Lean.
Its rate is $1/(2d4^d)$, instead of the earlier rate $2^{-O(d2^d)}$.
For an error margin $\varepsilon>0$, more than one half plus that margin forces
\[
 k<2d4^d\ln(2/\varepsilon).
\]
This improves the dependence on $\varepsilon$ from inverse-linear to logarithmic.

% New refinement. Root and path_stability independently audited all steps and constants.
\section{Conditioning on terminal data block by block}
\label{sec:second-order-block-conditioning}

A block is a maximal two-connected subgraph, with each bridge counted as a block.
An isolated vertex is a degenerate block when needed.
The cycle space of a graph is the direct sum of its block cycle spaces.
The block-cut graph of each connected component is a tree.
Its nodes are the blocks and the articulation vertices, with incidence as adjacency.

\begin{lemma}[Weighted conditional block estimate]\label{lem:weighted-block-conditioning}
Let $Q$ be a connected loopless multigraph, and let $S$ be a nonempty proper vertex set.
For each vertex $v$, let $L_v$ be its local edge-coordinate space in $\mathcal Z(Q)^*$.
Put
\[
 U=\sum_{s\in S}L_s,\qquad B=|\delta_Q(S)|.
\]
Let $\mathcal V\subseteq V(Q)\setminus S$ consist of vertices with $L_v\cap U\ne0$.
Write $I_v$ for the indicator that $v$ is closed, and put
$p_v=\mathbb P(I_v=1)=2^{-\dim L_v}$.
For arbitrary $0\le\eta_v\le1$, there is a set $E\subseteq\mathcal V$ with $|E|\le5B$ such that
\begin{equation}\label{eq:weighted-block-conditioning}
 \mathbb E\prod_{v\in\mathcal V}(1-\eta_v I_v)
 \le\frac12+\frac12\exp\!\left(-\sum_{v\in\mathcal V\setminus E}\eta_v p_v\right).
\end{equation}
\end{lemma}

\begin{proof}
Put $H=Q-S$.
In its block-cut forest, attach one distinct terminal leaf for each edge of $\delta_Q(S)$.
Attach that leaf to the articulation node of its endpoint in $H$, if that endpoint is an articulation.
Otherwise attach it to the unique block containing its endpoint.
Thus terminals at nonarticulation vertices are also retained.
Prune each tree to the union of paths between terminal leaves.
Components with fewer than two terminals contribute nothing.
Call the resulting forest $T$.
It has at most $B$ leaves.

Exclude the vertices of $\mathcal V$ adjacent to $S$, of which there are at most $B$.
The cut-space argument in Lemma~\ref{lem:local-terminal-intersection} applies to every remaining vertex.
It shows that such a vertex is an articulation separating at least two terminal-bearing components of $H-v$.
Its articulation node therefore has degree at least two in $T$.

Also exclude articulation nodes of degree at least three in $T$.
There are at most $B$ of these.
Finally exclude articulation nodes adjacent in $T$ to a block node of degree at least three.
The sum of the degrees of all branching nodes in a forest with at most $B$ leaves is at most $3B$.
This last exclusion therefore removes at most $3B$ further vertices.
Let $E$ be the union of the excluded vertices.

Each remaining vertex $v$ has exactly two terminal-bearing incident blocks in $H$.
All other incident blocks lie in branches containing no terminal.
The two terminal sides give the same nonzero functional $a_v\in U$.
The cut-space argument also shows
\begin{equation}\label{eq:block-intersection-line}
 L_v\cap U=\{0,a_v\}.
\end{equation}
Indeed, a local cut can only select a union of the components of $H-v$.
Terminal-free components contribute zero.
The two terminal-bearing components have complementary terminal sets, whose total functional is zero.

Join two remaining vertices when they belong to a common block of $H$.
Any shared block lies on the terminal forest between their articulation nodes.
It has degree two there, since neighbors of branching block nodes were excluded.
The resulting conflict graph is therefore a forest of paths.
Its two color classes contain no pair of vertices incident with the same block.
Choose a color class $J$ with
\[
 \sum_{v\in J}\eta_vp_v
 \ge\frac12\sum_{v\in\mathcal V\setminus E}\eta_vp_v.
\]

Condition on the evaluation of all functionals in $U$.
This determines every edge incident with $S$.
The remaining edge set in $H$ is a uniform member of an affine translate of $\mathcal Z(H)$.
Its random choices in different blocks are independent.

For clarity, zero-star feasibility can also be checked within each block.
At an articulation $v$, each incident block has a prescribed star parity, determined by the terminal data on its side.
The terminal-free sides have parity zero.
The other two sides have parity $a_v$.
Deleting $v$ from a nondegenerate block leaves it connected.
Hence, whenever the prescribed parity at $v$ is zero, a compatible edge set avoiding its star exists in that block.
This follows by joining its remaining odd vertices inside the connected graph obtained after deletion.
For a bridge block, parity zero simply forces its edge to be absent.
Thus the full star at $v$ can be zero exactly when $a_v=0$.

Equation~\eqref{eq:block-intersection-line} gives the exact conditional probability
\[
 \mathbb P(I_v=1\mid U)=2p_v\,\mathbf1_{\{a_v=0\}}.
\]
The indicators for $v\in J$ are conditionally independent, because their incident block sets are disjoint.
Ignoring the other factors and using $1-x\le e^{-x}$ gives
\[
 \mathbb E\!\left[\prod_{v\in\mathcal V}(1-\eta_vI_v)\,\middle|\,U\right]
 \le\exp\!\left(-2\sum_{v\in J}\eta_vp_v\mathbf1_{\{a_v=0\}}\right).
\]

Put $W=\sum_{v\in J}\eta_vp_v\mathbf1_{\{a_v=0\}}$ and $A=\sum_{v\in J}\eta_vp_v$.
Each $a_v$ is a nonzero functional, so $\mathbb EW=A/2$, while $0\le W\le A$.
If $A=0$, the claimed bound is trivial.
Otherwise the chord bound for the convex function $e^{-2x}$ gives
\[
 \mathbb E e^{-2W}\le\frac12+\frac12e^{-2A}.
\]
The choice of $J$ now proves~\eqref{eq:weighted-block-conditioning}.
No independence among the functionals $a_v$ is required.
\end{proof}

\begin{theorem}[A sharper exponential cut constant]\label{thm:sharper-exponential-cuts}
Let $Q$ be a connected loopless multigraph with $k\ge1$ marked vertices of degree at most $d\ge2$.
For a uniform even subgraph, let $Z$ count the closed marked vertices.
Then
\begin{equation}\label{eq:sharper-exponential-cuts}
 \mathbb E2^{-Z}\le\frac12+2\exp\!\left(-\frac{k}{8d2^d}\right).
\end{equation}
The same bound holds for simultaneous forest restrictions to $k$ disjoint connected cyclic regions with boundaries of size at most $d$.
\end{theorem}

\begin{proof}
A zero local space gives an expectation at most one half.
A direct sum of $2^d$ local spaces gives an expectation below one half, as before.
Otherwise choose a maximal direct-sum family with vertex set $S$ and size $t<2^d$.
Every other marked local space meets $U=\sum_{s\in S}L_s$ nontrivially.

Suppose first that $k\ge2(5d+1)2^d$.
Then $S$ is nonempty and proper.
Its boundary has size $B\le dt$.
Apply Lemma~\ref{lem:weighted-block-conditioning} with $\eta_v=1/2$ to the marked vertices outside $S$.
At least $k-t-5B$ marked vertices remain outside its exceptional set.
Each has $p_v\ge2^{1-d}$, so
\[
 \mathbb E2^{-Z}-\frac12
 \le\frac12\exp\!\left(-2^{-d}(k-t-5B)\right)
 \le\frac12\exp\!\left((5d+1)-\frac{k}{2^d}\right)
 \le\frac12\exp\!\left(-\frac{k}{2^{d+1}}\right).
\]
This is stronger than~\eqref{eq:sharper-exponential-cuts}.

If $k<2(5d+1)2^d$, then
\[
 \frac{k}{8d2^d}<\frac{5d+1}{4d}\le\frac{11}{8}<\ln4.
\]
The right side of~\eqref{eq:sharper-exponential-cuts} therefore exceeds one.
This proves the bound for all $k$.
The cyclic-region statement follows from the same quotient and internal-cycle-space argument as before.
\end{proof}

\paragraph{Weighted scope.}
Lemma~\ref{lem:weighted-block-conditioning} retains the exact closure probabilities $p_v$.
It loses only the weighted mass on its exceptional vertices and one half of the remaining mass in the path coloring.
Its bound depends on no maximum degree after the exceptional set has been chosen.
This identifies a possible route to a weighted estimate for regions with unequal boundaries.

% Source proofs adapted from pancyclic-paper.tex, Section6.
% Original source is unchanged. This file makes the new draft self-contained.
\section{Cycle-space, walk, and spectral estimates}
\label{sec:polynomial-foundations}

A region is a nonempty connected vertex set.
For a region $U$, let $\delta U$ be its edge boundary.
Write $\mathcal O_U$ for the event that a uniform even subgraph avoids $\delta U$.
For disjoint regions $U,V$, their links are the edges between them and the components of $G-U-V$ adjacent to both.
The number of links is $\ell(U,V)$.
The following proofs are adapted from Section6 of the existing manuscript.
They are included to expose every dependency of the polynomial argument.

\begin{lemma}[Link formula]\label{lem:link}
Let $A$ and $B$ be disjoint regions of a connected graph $G$. Then
$\PP(\cO_A)=2^{\comp(G-A)-|\delta A|}\ge2^{-|\delta A|}$ and
\[
  \PP(\cO_A\cap\cO_B)=2^{\ell(A,B)-1}\,\PP(\cO_A)\,\PP(\cO_B).
\]
Conditionally on $\cO_A$, the set $\xi\cap E(G[A])$ is uniform on $\Z(G[A])$; conditionally on
$\cO_A\cap\cO_B$, the sets $\xi\cap E(G[A])$ and $\xi\cap E(G[B])$ are independent and uniform on
$\Z(G[A])$ and $\Z(G[B])$.
\end{lemma}

\begin{proof}
The even subgraphs that avoid $\delta A$ are the sums of an even subgraph of $G[A]$ and an even
subgraph of $G-A$. So $\PP(\cO_A)=2^{\beta(G[A])+\beta(G-A)-\beta(G)}$, and counting edges and
vertices shows that the exponent is $\comp(G-A)-|\delta A|$. In the same way
\[
  \PP(\cO_A\cap\cO_B)=2^{\beta(G[A])+\beta(G[B])+\beta(G-A-B)-\beta(G)}
     =2^{\comp(G-A-B)+1-|\delta A|-|\delta B|+e(A,B)},
\]
where $e(A,B)$ is the number of edges between $A$ and $B$. Let $a$, $b$ and $k$ be the numbers of
components of $G-A-B$ that are joined to $A$ only, to $B$ only, and to both. Then
$\comp(G-A-B)=a+b+k$, while $\comp(G-A)=a+1$ and $\comp(G-B)=b+1$, and the formula follows since
$\ell(A,B)=k+e(A,B)$. The conditional laws are read off from the two direct sum decompositions.
\end{proof}

So two cut events are independent when there is one link. They are correlated by a factor of
exactly two when there are two links. The correlation is stronger only for pairs with at least
three links, and the next lemma says that such pairs are rare.

\begin{lemma}[Exceptional partners]\label{lem:exc}
Let $A$ be a region of a connected graph $G$, and let $B_1,\dots,B_k$ be pairwise disjoint regions,
disjoint from $A$, with $\ell(A,B_i)\ge3$ for all $i$. Then $k\le|\delta A|$.
\end{lemma}

\begin{proof}
Contract $A$ to a vertex $a$ and each $B_i$ to a vertex $b_i$, and delete the loops that arise.
This changes neither $d=|\delta A|$ nor the numbers $\ell(A,B_i)$, because the sets $B_j$ with
$j\ne i$ are connected and are contracted inside components of $G-(A\cup B_i)$. Let $F$ be a spanning forest of
$G-a$, and for every edge $e=au$ at $a$ attach a new leaf $\tau_e$ to $u$. Call the resulting forest
$F^{+}$ and its $d$ new leaves \emph{terminals}.

Fix $i$. Every link between $a$ and $b_i$ yields a path of $F^{+}$ from $b_i$ to a terminal. For an
edge $e=ab_i$ it is the edge $b_i\tau_e$. For a component $Y$ of $G-\{a,b_i\}$ joined to both, choose
an edge $e=au$ with $u\in Y$ and take the path of $F^{+}$ from $b_i$ to $\tau_e$; after its first
vertex it stays in $Y\cup\{\tau_e\}$. Paths that come from different links therefore start with
different edges at $b_i$, and the union of any two of them is a path between two terminals. Let
$F^{\circ}\subseteq F^{+}$ be the union of all paths between terminals. As $\ell(a,b_i)\ge3$, the
vertex $b_i$ has degree at least $3$ in $F^{\circ}$. All leaves of $F^{\circ}$ are terminals, and a
forest with at most $d$ leaves has at most $d-2$ vertices of degree at least $3$. Hence $k\le d-2$ if $k\ge1$.
\end{proof}


\subsection{Walks and normalized traces}
Let $J$ be a simple graph of maximum degree three and girth greater than $D\ge3$.
Put $s=\lfloor(D-1)/2\rfloor$ and let $L>D$.
Let $N_\ell(u,w)$ count non-backtracking walks of length $\ell$ from $u$ to $w$.
Write $v=|V(J)|$ and $d_J=\sum_u(3-\deg_Ju)$.
For the spectral supply estimate below, assume $d_J/v<1$.

\begin{lemma}\label{lem:walks}
If $\ell>s$, then $N_\ell(u,w)\le3\cdot2^{\ell-s-1}$ for all $u,w$. Consequently, for every vertex
$u$ of $J$, the cycles $C$ of $J$ through $u$ with $|C|\le L$ satisfy $\sum_C2^{-|C|}\le L\,2^{-s}$.
\end{lemma}

\begin{proof}
There are at most $3\cdot2^{\ell-s-1}$ ways to choose the first $\ell-s$ steps. The last $s$ steps
form a non-backtracking walk of length $s$ between two given vertices. Such a walk is a path,
because $s$ is less than the girth, and there is at most one such path, because two of them would
give a cycle of length at most $2s<D$. For the second statement, a cycle of length $\ell$ through $u$
can be traversed from $u$ in two directions, so there are at most $\half N_\ell(u,u)$ of them, and
$\ell>D>s$.
\end{proof}

Let $B$ be the non-backtracking matrix of $J$. Its rows and columns are indexed by the directed
edges $(i\to j)$ of $J$, two for each edge. The entry in row $(i\to j)$ and column $(j'\to k)$ is
$1$ if $j'=j$ and $k\ne i$, and it is $0$ otherwise.

\begin{lemma}\label{lem:spectrum}
For every even $\ell\ge2$,
\[
  \tr B^{\ell}\;\ge\;\bigl(2-d_J/v\bigr)^{\ell}-3v\,2^{\ell/2}.
\]
\end{lemma}

\begin{proof}
Let $A$ be the adjacency matrix of $J$ and $\Delta$ the diagonal matrix of its degrees, and put
$M(\lambda)=\lambda^{2}I-\lambda A+\Delta-I$.

(1) \emph{If $M(\lambda)g=0$ for a real $\lambda>1$ and a vector $g\ne0$, then $\lambda$ is an
eigenvalue of $B$.} A direct computation shows that $f(i\to j)=\lambda g(j)-g(i)$ satisfies
$Bf=\lambda f$; and $f=0$ would give $g(i)=\lambda g(j)=\lambda^{2}g(i)$ along every edge, hence
$g=0$ (at an isolated vertex, $M(\lambda)g=0$ already says $(\lambda^{2}-1)g=0$).

(2) \emph{If $Bf=\lambda f$ with $f\ne0$ and $\lambda^{2}\ne1$, then $M(\lambda)g=0$ for
$g(j)=\sum_{k\sim j}f(j\to k)$, and $g\ne0$.} The eigenvalue equation at $(i\to j)$ and at
$(j\to i)$ reads $\lambda f(i\to j)=g(j)-f(j\to i)$ and $\lambda f(j\to i)=g(i)-f(i\to j)$.
Eliminating $f(j\to i)$ gives $(\lambda^{2}-1)f(i\to j)=\lambda g(j)-g(i)$. This shows that
$g\ne0$. Summing it over the neighbours $j$ of $i$ gives $M(\lambda)g=0$.

(3) \emph{$B$ has a real eigenvalue $\lambda_1\ge2-d_J/v$.} The matrix
$M(2)=2(\Delta-A)+(3I-\Delta)$ is positive semidefinite, while
$\mathbf{1}^{\!\top}M(\lambda)\mathbf{1}=(\lambda-1)(v\lambda-2v+d_J)$ vanishes at
$\lambda^{*}=2-d_J/v>1$. The smallest eigenvalue of the symmetric matrix $M(\lambda)$ is continuous
in $\lambda$, at most $0$ at $\lambda^{*}$ and at least $0$ at $2$. So $M(\lambda_1)$ is singular for
some $\lambda_1\in[\lambda^{*},2]$, and (1) applies.

(4) \emph{Every non-real eigenvalue $\lambda$ of $B$ satisfies $|\lambda|\le\sqrt2$.} With $g$ as
in~(2), $g^{*}M(\lambda)g=0$ is a quadratic equation $a\lambda^{2}-b\lambda+c'=0$ with real
coefficients $a=g^{*}g>0$, $b=g^{*}Ag$ and $c'=g^{*}(\Delta-I)g\le2a$. Its roots are $\lambda$ and
$\bar\lambda$, so $|\lambda|^{2}=c'/a\le2$.

Now $\tr B^{\ell}$ is the sum of the $\ell$th powers of the $2|E(J)|\le3v$ eigenvalues of $B$. For
even $\ell$ the real eigenvalues contribute nonnegative terms, one of which is at least
$(2-d_J/v)^{\ell}$, and the non-real ones contribute at least $-3v\,2^{\ell/2}$ in total.
\end{proof}

The argument uses the same matrix-pencil proof as the existing manuscript.
It needs no expansion assumption.


\begin{lemma}[Repeated closed walks]\label{lem:polynomial-repeated-walks}
Let $\mathcal C$ contain the cycles of $J$ with length at most $L$ and put $\mu=\sum_{C\in\mathcal C}2^{-|C|}$.
Then
\[
 \mu\le T:=\sum_{\ell\le L}\frac{\operatorname{tr}B^\ell}{2\ell2^\ell}
 \le\left(1+\tfrac32L^2 2^{-s}\right)\mu.
\]
\end{lemma}
\begin{proof}
$\tr B^{\ell}$ is the number of closed walks $u_0u_1\cdots u_\ell=u_0$ in $J$ with
$u_{i+1}\ne u_{i-1}$ for all $i$ modulo $\ell$. A cycle of length $\ell$ gives $2\ell$ of them; let
$n_\ell$ be the number of the others. Then
\[
  T:=\sum_{\ell\le L}\frac{\tr B^{\ell}}{2\ell\,2^{\ell}}=\mu+\sum_{\ell\le L}\frac{n_\ell}{2\ell\,2^{\ell}} .
\]
Consider a walk counted by $n_\ell$, and let $j$ be the least index with
$u_j\in\{u_0,\dots,u_{j-1}\}$, say $u_j=u_i$. Then $j<\ell$, the walk $C=u_iu_{i+1}\cdots u_j$ is a
cycle, of length $q$ say, and $u_ju_{j+1}\cdots u_\ell u_1\cdots u_i$ is a non-backtracking closed
walk of length $\ell-q>0$ from $u_i$ to itself. It contains a cycle, so $\ell-q>D>s$. The original
walk is determined by $C$, by the vertex $u_i$ and a direction on $C$ ($2q$ choices), by the
complementary walk (at most $3\cdot2^{\ell-q-s-1}$ choices, by Lemma~\ref{lem:walks}), and by the
position of $u_0$ on the latter (at most $\ell$ choices). Hence
\[
  \sum_{\ell\le L}\frac{n_\ell}{2\ell\,2^{\ell}}
  \;\le\;\sum_{\ell\le L}\ \sum_{C\in\mathcal{C}}\tfrac32\,|C|\,2^{-|C|-s}
  \;\le\;\tfrac32\,L^{2}\,2^{-s}\,\mu,
\]
as required.


\end{proof}

In particular, suppose $v=O(r)$, $x=\log_2r$, $L=\lfloor Ax\rfloor$, and $A>6$ is fixed.
If $Ld_J=o(v)$ and $L^2 2^{-s}=o(1)$, then
\[
 \mu\ge(1/4-o(1))\ln(A/6).
\]
Indeed, sum the even trace bound over $6x\le\ell\le L$.
Its normalized leading term is $1-o(1)$, and its nonreal error is $O(r^{-2})$.
The even harmonic sum is $(1/4-o(1))\ln(A/6)$.
The repeated-walk error is $o(\mu)$.

\subsection{The uniform bound on links}
Suppose $\Gamma$ is connected and subcubic, $\Gamma[P]$ is connected, and every vertex outside $P$ has degree three.
Let $\mathcal Q$ be a maximal disjoint family of cycles of length at most $D$ in $\Gamma-P$.
Write $S$ for its vertex union and $J=\Gamma-(P\cup S)$.
Assume the probability of a forest outside $P$ exceeds $1/2+\varepsilon$.
For disjoint cycles $C,C'$ of $J$ with lengths at most $L$, and an integer $t$ with $2t+D/2\le D$,
\begin{equation}\label{eq:polynomial-foundation-links}
 \ell(C,C')\le1+\frac{2L}{t}+8t2^t\ln(2/\varepsilon)+(8L/D)^2.
\end{equation}

To prove this, let $U=V(C)\cup V(C')$.
At most one component of $\Gamma-U$ meets $P$.
At most $2L/t$ other components have more than $t$ boundary edges.
The exponential cut estimate bounds the number of cyclic components with at most $t$ boundary edges by $8t2^t\ln(2/\varepsilon)$.

A remaining tree component $Y$ has $|Y|=|\delta Y|-2\le t-2$.
It contains no chosen short cycle, since every chosen cycle avoids $U$ and is connected.
Thus $Y\subseteq J$.
If $Y$ is adjacent to both cycles, it supplies a connecting path of length at most $t$.
Direct edges between the cycles are also such paths.
Subcubic degree makes these paths vertex-disjoint, including their endpoints.
Divide each cycle into at most $8L/D$ arcs with at most $D/4$ consecutive vertices.
Two connecting paths with the same pair of endpoint arcs would produce a cycle of length at most $2t+D/2\le D$ in $J$.
There is therefore at most one path per pair of arcs.
This proves \eqref{eq:polynomial-foundation-links}.

\subsection{Component variables and exceptional pairs}
For a cycle $C\subseteq J$, put $a_C=2^{-|C|}$ and let $I_C$ mean that $C$ is a component of the uniform even subgraph $\xi$ of $\Gamma$.
Set $p_C=\Pr(I_C)$ and $X_C=(a_C/p_C)\mathbf1_{I_C}$.
The condition $E(C)\subseteq\xi$ is feasible and imposes at most $|C|$ binary linear equations.
Hence $p_C\ge a_C$, so $\mathbb EX_C=a_C$ and $\mathbb EX_C^2\le a_C$.
The link formula gives
\[
 \mathbb E(X_CX_{C'})=a_Ca_{C'}2^{\ell(C,C')-1}
\]
for disjoint cycles.
Distinct intersecting cycles cannot both be components.

Fix $C$ and choose a maximal disjoint family of partners with at least three links to $C$.
The exceptional-partner lemma bounds its size by $L$.
Its vertex union has size at most $L^2$ and meets every exceptional partner.
The local walk bound therefore bounds their total weight by $L^3 2^{-s}$.
If all link counts are at most $b$, their total normalized joint contribution is at most
\[
 a_C2^bL^3 2^{-s}.
\]
These estimates apply to each of the two graph deletions in the next section.

\section{A polynomial unrestricted error from two girth scales}
\label{sec:two-scale-polynomial-error}

This proof combines Fable's nested-cycle argument with the exponential cut estimate.
It uses two nested short-cycle deletions.
The larger deletion finds a long chain.
Restoring the medium-length deleted cycles makes most chain blocks suitable for a spectral test.
The proof is mathematical and has not been formalized in Lean.

All logarithms in this section have base two unless written as $\ln$.
Fix
\[
 A=6e^{168},\qquad t=\lceil10000A\rceil,
 \qquad \kappa=\frac1{10^6t2^t}>0.
\]
All thresholds below are absolute.
They are not optimized.

\begin{theorem}[A polynomial forest error]
\label{thm:two-scale-polynomial-forest}
Let $\Gamma$ be a connected multigraph of rank $r$ and maximum degree at most three.
Let $P$ be a nonempty vertex set such that $\Gamma[P]$ is connected,
\[
 |P|\le\frac{2r}{\log^2r},
\]
and every vertex outside $P$ has degree three.
For a uniform even subgraph $\xi$, let $\mathcal F$ mean that $\xi$ is a forest on $\Gamma-P$.
For all sufficiently large $r$,
\[
 \Pr(\mathcal F)\le\frac12+r^{-\kappa}.
\]
\end{theorem}

\subsection{The two cycle families and their moments}

Assume, for a contradiction, that
\begin{equation}\label{eq:two-scale-assumption}
 \Pr(\mathcal F)>\frac12+\varepsilon,
 \qquad \varepsilon=r^{-\kappa}.
\end{equation}
Put
\[
 x=\log r,\qquad L=\lfloor Ax\rfloor,
 \qquad D_0=\lfloor x/100\rfloor,
 \qquad D_1=\lfloor x/10\rfloor,
 \qquad s_j=\lfloor(D_j-1)/2\rfloor.
\]
First choose a maximal disjoint family $\mathcal Q_0$ of cycles of length at most $D_0$ in $\Gamma-P$.
Extend this family to a maximal disjoint family $\mathcal Q_1$ of cycles of length at most $D_1$.
Let $S_j$ be the union of the vertices of $\mathcal Q_j$, and put
\[
 J_j=\Gamma-(P\cup S_j).
\]
Thus $S_0\subseteq S_1$ and $J_1\subseteq J_0$.

Every chosen cycle is a cyclic region with boundary at most its length.
Theorem~\ref{thm:sharper-exponential-cuts} and \eqref{eq:two-scale-assumption} give
\[
 |\mathcal Q_j|<8D_j2^{D_j}\ln(2/\varepsilon).
\]
Consequently,
\begin{equation}\label{eq:two-scale-deletion-sizes}
 |S_0|=O(x^3r^{1/100}),\qquad
 |S_1|=O(x^3r^{1/10}).
\end{equation}
Each $J_j$ is simple and has girth greater than $D_j$.
Its order $v_j$ and degree deficit $d_j$ satisfy
\[
 v_j=2r+O(|P|+|S_j|),\qquad
 d_j\le3(|P|+|S_j|),\qquad Ld_j=o(v_j).
\]

Let $\mathcal C_j$ be the cycles of $J_j$ of length at most $L$.
Write
\[
 a_C=2^{-|C|},\qquad \mu_j=\sum_{C\in\mathcal C_j}a_C.
\]
The manuscript's walk bound gives, uniformly over vertices $v$ of $J_j$,
\begin{equation}\label{eq:two-scale-local-mass}
 \sum_{C\in\mathcal C_j:v\in C}a_C\le L2^{-s_j}.
\end{equation}
The spectral supply argument, applied to $J_1$, gives
\begin{equation}\label{eq:two-scale-initial-mass}
 \mu_1\ge\left(\frac18-o(1)\right)\ln(A/6)>20.
\end{equation}
Here $L^2 2^{-s_1}=o(1)$ controls repeated non-backtracking walks, and $Ld_1=o(v_1)$ controls the spectral loss.

We record the moment bounds at both girth scales.
For disjoint $C,C'\in\mathcal C_j$, the manuscript's link decomposition now gives
\[
 \ell(C,C')\le1+\frac{2L}{t}
   +8t2^t\ln(2/\varepsilon)+(8L/D_j)^2.
\]
The first term uses the connectedness of $P$.
The third term uses the exponential cut estimate for the cyclic complementary regions.
The tree-region argument remains valid because $2t+D_j/2\le D_j$ for large $r$.
The choice of constants gives
\begin{equation}\label{eq:two-scale-link-bound}
 \ell(C,C')\le x/1000
\end{equation}
for large $r$, at both scales.
Indeed, $2L/t\le x/5000$, the third term equals
$(8\ln2/10^6)x+O_t(1)$, and the last term is bounded independently of $r$.

Use the manuscript's normalized cycle-component variables
\[
 X_C=\frac{a_C}{\Pr(C\text{ is a component of }\xi)}
       \mathbf1_{\{C\text{ is a component of }\xi\}},
 \qquad Z_j=\sum_{C\in\mathcal C_j}X_C.
\]
They satisfy $\mathbb EX_C=a_C$, $\mathbb EX_C^2\le a_C$, and
\[
 \mathbb E(X_CX_{C'})=a_Ca_{C'}2^{\ell(C,C')-1}
\]
for disjoint cycles.
Distinct intersecting cycles have zero joint contribution.
The exceptional-partner lemma and \eqref{eq:two-scale-link-bound} bound the contribution of partners with at least three links by
\[
 a_C\,2^{x/1000}L^3 2^{-s_j}=o(a_C).
\]
This holds even at the smaller scale, since $s_0=x/200+O(1)$.

Let $\Lambda_j$ be the total weight of ordered disjoint one-link pairs in $\mathcal C_j$.
There is a quantity $\eta_r\to0$ such that
\begin{equation}\label{eq:two-scale-inverse-moment}
 \mathbb EZ_j^2
 \le2\mu_j^2-\Lambda_j+(1+\eta_r)\mu_j.
\end{equation}
On $\mathcal F$, one has $Z_j=0$.
Cauchy--Schwarz and \eqref{eq:two-scale-assumption} therefore imply
\begin{equation}\label{eq:two-scale-two-consequences}
 \Lambda_j<(1+\eta_r)\mu_j,
 \qquad
 \mu_j\le\frac{1+\eta_r}{4\varepsilon}.
\end{equation}

\subsection{A chain exists after weighted averaging}

For disjoint cycles $C,C'\in\mathcal C_1$, call $C'$ \emph{far from $P$ behind $C$}
when it lies in a component of $\Gamma-V(C)$ that avoids $P$.
This relation is antisymmetric.
Indeed, a path from $P$ to its first vertex of $C$ avoids that far component.
It therefore avoids $C'$.
Since $C$ is connected and disjoint from $C'$, all of $C$ lies in the $P$-component of $\Gamma-V(C')$.
It follows that the total weight of ordered far pairs is at most $\mu_1^2/2$.

By \eqref{eq:two-scale-local-mass}, the ordered weight of intersecting pairs is at most
$\mu_1L^2 2^{-s_1}=o(\mu_1)$.
The same bound applies to disjoint pairs joined by an edge.
By \eqref{eq:two-scale-two-consequences}, one-link pairs have weight at most $(1+o(1))\mu_1$.
Thus the ordered weight of the remaining pairs is at least
\[
 \frac12\mu_1^2-(1+o(1))\mu_1.
\]
In every remaining pair, $C'$ lies in the $P$-component of $\Gamma-V(C)$,
no edge joins the two cycles, and they have at least two links.
Equation~\eqref{eq:two-scale-initial-mass} therefore supplies one cycle $C$ whose such partners have total weight at least eight.

Put $H=\Gamma-V(C)$ and let $N(C)$ be the vertices outside $C$ adjacent to $C$.
Since $C$ lies outside $P$, every vertex of $C$ is cubic, so $|N(C)|\le|C|\le L$.
For each selected partner $C'$, at least two components of $H-V(C')$ are adjacent to both cycles.
At most one contains $P$.
Choose another such component $Y$ and a vertex $u\in Y\cap N(C)$.
The component $Y$ lies in the main component of $H$, since it is adjacent to $C'$ there.
Thus $C'$ separates $u$ from $P$ in $H$.

For some fixed $u\in N(C)$, cycles separating $u$ from $P$ in $H$ have total weight at least $8/L$.
Choose a maximal vertex-disjoint subfamily of these cycles.
If it has $T$ members, its vertex union meets every cycle in the family.
Equation~\eqref{eq:two-scale-local-mass} gives
\[
 \frac8L\le TL^2 2^{-s_1},\qquad
 T\ge\frac{8\,2^{s_1}}{L^3}
 \ge\frac{r^{1/20}}{L^3}
\]
for large $r$.

These disjoint cycles are nested between $u$ and $P$ in $H$.
Here is the required separation argument.
For two such cycles, replace the middle of a $u$--$P$ path by an arc of the first cycle.
The resulting path meets the second cycle in either the $u$-component or the $P$-component after deleting the first.
Connectedness places the entire second cycle in that component.
The two cycles cannot each lie on the other's $u$-side, since a path through the first cycle would then avoid the second separator.
They cannot each lie on the other's $P$-side, by the reversed argument.
Thus they have a consistent linear order.

Write the ordered cycles as $C_1,\ldots,C_T$.
Let $U_i$ be the component of $u$ in $H-V(C_i)$.
Then $U_i\subseteq U_{i+1}$, and $C_i\subseteq U_{i+1}$.

\subsection{Restore medium cycles inside the chain}

Put $q=\lceil20000Ax\rceil+2$.
Choose disjoint index intervals $[a,b]$ with $b-a=q$, leaving a separator cycle between consecutive intervals.
There are
\[
 m_* = \Omega\!\left(\frac{r^{1/20}}{x^4}\right)
\]
such intervals.
For each interval define the full vertex set
\[
 V(a,b)=U_b\setminus\bigl(U_a\cup V(C_a)\bigr).
\]
It contains $C_{a+1},\ldots,C_{b-1}$.
Consequently,
\begin{equation}\label{eq:two-scale-block-volume}
 |V(a,b)|\ge(q-1)(D_1+1)\ge2000Ax^2.
\end{equation}
These sets are pairwise disjoint and pairwise non-adjacent in $H$.
Every edge of $H$ leaving one set ends on $C_a$ or $C_b$.
Its boundary in $H$ is therefore at most $2L$.
The set has no neighbor in $P$, since $C_b$ separates it from $P$.

The definition uses full vertex sets in $H$.
It restores the vertices in $S_1\setminus S_0$ that lie between the separator cycles.
This restoration is the reason for using two girth scales.

Discard every block meeting $N(C)$.
At most $L$ blocks are lost.
Also discard every block meeting $S_0$ or a neighbor of $S_0$ in $\Gamma$.
At most $4|S_0|$ further blocks are lost, since $\Gamma$ has maximum degree three.
By \eqref{eq:two-scale-deletion-sizes}, this loss is $o(m_*)$.
Let $V_1,\ldots,V_m$ be the remaining blocks.
Then
\begin{equation}\label{eq:two-scale-many-blocks}
 m=\Omega\!\left(\frac{r^{1/20}}{x^4}\right).
\end{equation}

Each $V_i$ lies in $J_0$.
It has no neighbor in $P\cup S_0$, so every one of its vertices has degree three in $J_0$.
Thus its degree deficit in $J_0$ is exactly zero.
Its boundary in $J_0$ has size at most $2L$.
The blocks remain pairwise non-adjacent.

\subsection{Many blocks give many spectral modes}

We include the spectral argument for completeness.
Let $J$ be a simple graph of maximum degree three.
Suppose $V_1,\ldots,V_m$ are pairwise disjoint and pairwise non-adjacent vertex sets.
Write $b_i=|\delta_J V_i|$ and
$\delta_i=\sum_{v\in V_i}(3-\deg_J v)$.
If $0<\theta<1/2$ and
\[
 \delta_i+3b_i\le\theta|V_i|,
\]
then the non-backtracking matrix of $J$ has at least $m$ real eigenvalues in $[2-\theta,2]$, with multiplicity.

To prove this, let $A_J$ be its adjacency matrix and $\Delta_J$ its degree matrix.
Put $M(\lambda)=\lambda^2I-\lambda A_J+\Delta_J-I$.
For the indicator vector of $V_i$,
\[
 \mathbf1_{V_i}^{\mathsf T}M(2-\theta)\mathbf1_{V_i}
 =-\theta(1-\theta)|V_i|
 +(1-\theta)\delta_i+(2-\theta)b_i\le0.
\]
Nonadjacency makes the corresponding cross terms zero.
Hence $M(2-\theta)$ has at least $m$ nonpositive eigenvalues.
The matrix $M(2)$ is positive semidefinite.
Its ordered eigenvalues vary continuously, so at least $m$ zeros occur in $[2-\theta,2]$, counted with kernel multiplicity.
Equivalently, $M'(\lambda)=2\lambda I-A_J$ is positive definite for $\lambda>3/2$, so each ordered eigenvalue crosses zero at most once there.
The manuscript's map
$g\mapsto f$, with $f(v\to w)=\lambda g(w)-g(v)$,
injects each kernel into the corresponding non-backtracking eigenspace.
This proves the eigenvalue count.

For even $\ell$, the nonreal eigenvalues have modulus at most $\sqrt2$.
Therefore
\[
 \operatorname{tr}B_J^\ell
 \ge m(2-\theta)^\ell-3|V(J)|2^{\ell/2}.
\]
Apply this to the clean blocks in $J_0$, with $\theta=1/(12x)$.
Their zero deficit, boundary at most $2L$, and \eqref{eq:two-scale-block-volume} verify the hypothesis.
Sum the normalized trace over even lengths $6x\le\ell\le24x$.
The first term contributes at least a fixed positive constant times $m$.
The nonreal error is $o(1)$, since $|V(J_0)|=O(r)$.
The repeated-walk bound is at most
$\tfrac32L^2 2^{-s_0}\mu_0=o(\mu_0)$.
The manuscript's trace-to-cycle calculation consequently gives, for large $r$,
\[
 \mu_0\ge m/16-1.
\]
Together with \eqref{eq:two-scale-many-blocks}, this yields
\[
 \mu_0=\Omega\!\left(\frac{r^{1/20}}{x^4}\right).
\]
But \eqref{eq:two-scale-two-consequences} gives
\[
 \mu_0\le\frac{1+o(1)}4\,r^\kappa.
\]
Since $\kappa<1/100$, these bounds contradict each other for large $r$.
This proves Theorem~\ref{thm:two-scale-polynomial-forest}.

\subsection{Transfer to unrestricted cycle counts}

\begin{corollary}[A polynomial unrestricted cycle error]
\label{cor:two-scale-polynomial-cycles}
For the absolute constant $\kappa$ above,
\[
 \psi(r)\le2^{r-1}+(1+o(1))\frac{2^r}{r^\kappa}.
\]
In particular, $\psi(r)-2^{r-1}=O(2^r/r^\kappa)$.
\end{corollary}

\begin{proof}
Normalize a graph by deleting irrelevant trees and suppressing degree-two vertices.
Rank and cycle count are preserved.
The resulting kernel has at most $2r-2$ vertices and a cycle of length $O(\log r)$.
If necessary, subdivide its edges uniformly to work with a simple graph.
The chosen cycle still has length $O(\log r)$.

For this connected cycle witness, the manuscript's marked cubic reduction has connected marked set $P$ and
$|P|=O(\log r)\le2r/\log^2r$.
Theorem~\ref{thm:two-scale-polynomial-forest} therefore bounds the original forest probability by $1/2+r^{-\kappa}$.
The cycle-flip injection for any witness edge gives
\[
 c_e(G)\le2^{r-2}+2^{r-1}r^{-\kappa}.
\]
Deleting that cyclic edge preserves connectedness and reduces rank by one.
The established edge recurrence, with $E_r=\psi(r)-2^{r-1}$, gives
\[
 E_r\le\max\{0,E_{r-1}+2^{r-1}r^{-\kappa}\}
\]
for all sufficiently large $r$.
Iterating and summing the geometric tail yields
\[
 E_r\le O(1)+\sum_{j\le r}2^{j-1}j^{-\kappa}
      =(1+o(1))2^r r^{-\kappa}.
\]
This proves the corollary.
\end{proof}

\section{The critical marked-set scale and a direct cycle transfer}
\label{sec:two-scale-extensions}

This section records extensions separately from the audited two-scale proof.
Retain its absolute constants $A,t,\kappa$.

\begin{theorem}[Several marked components at the critical scale]
\label{thm:two-scale-critical-marking}
Let $\Gamma$ be a connected multigraph of rank $r$ and maximum degree at most three.
Let $P$ be a nonempty vertex set containing every vertex of degree less than three.
Suppose
\[
 |P|\le\frac{r}{1000A\log r},\qquad
 c:=\operatorname{comp}(\Gamma[P])\le\frac{\log r}{10000}.
\]
For a uniform even subgraph $\xi$, let $\mathcal F$ mean that its restriction to $\Gamma-P$ is a forest.
For all sufficiently large $r$,
\[
 \Pr(\mathcal F)\le\frac12+r^{-\kappa}.
\]
\end{theorem}

\begin{proof}
Use the two nested short-cycle deletions and the same parameters as in Theorem~\ref{thm:two-scale-polynomial-forest}.
We verify the three changes required by the new hypotheses.

First, the marked-set deficit causes a fixed small spectral loss.
For $J_1$, one has $v_1=2r+o(r)$ and
\[
 d_1\le\frac{3r}{1000Ax}+o(r/x),\qquad x=\log r.
\]
Consequently,
\[
 \frac{Ld_1}{2v_1}\le\frac3{4000}+o(1)<\frac1{500}.
\]
For each even length $6x\le\ell\le L$, the normalized Perron contribution is at least $1-1/500$.
The even harmonic sum tends to $\tfrac14\ln(A/6)=42$.
The repeated-walk error remains $o(\mu_1)$.
Thus $\mu_1>20$ for large $r$.
This uses a fixed spectral margin, rather than a vanishing spectral loss.

Second, the link bound has $c$ in place of one.
Its other terms are unchanged, and
\[
 c+\frac{2L}{t}+8t2^t\ln(2r^\kappa)+(8L/D_0)^2
 \le x/1000
\]
for large $r$.
Both inverse moment bounds therefore remain valid.

Third, fix one connected component $P_0$ of $\Gamma[P]$.
Use $P_0$ as the anchor in the definitions of main and far components.
The far relation remains antisymmetric, since $P_0$ is connected.
The same weighted averaging gives a long chain of cycles separating one vertex $u$ from $P_0$.
Every cycle in this chain avoids all of $P$.

Form the full blocks $V(a,b)$ in $\Gamma-V(C)$, as before.
Each connected component of $\Gamma[P]$ can meet at most one block.
Indeed, each block boundary ends on its two separator cycles, both disjoint from $P$.
A connected set inside $P$ cannot cross such a boundary.
Discard every block meeting $P$, losing at most $c$ blocks.
Every retained block has no neighbor in $P$ for the same boundary reason.

Now discard blocks meeting $N(C)$ and blocks meeting $S_0$ or a neighbor of $S_0$.
The additional loss is at most $L+4|S_0|$.
The total loss is still negligible compared with $\Omega(r^{1/20}/x^4)$.
Every retained block has zero deficit in $J_0$, boundary at most $2L$, and the previous volume lower bound.
The multi-mode trace argument again contradicts
$\mu_0\le(1+o(1))r^\kappa/4$.
\end{proof}

The coefficient $1/(1000A)$ is deliberately small.
This theorem reaches the order $r/\log r$ with a polynomial error above one half.
It does not assert the same bound for every fixed coefficient at that scale.

\begin{proposition}[Direct transfer through a cubic extremizer]
\label{prop:two-scale-direct-cycle-transfer}
The polynomial forest theorem alone implies
\[
 \psi(r)\le2^{r-1}+(1+o(1))2^r r^{-\kappa}.
\]
This transfer needs no marked cubic reduction.
\end{proposition}

\begin{proof}
Choose a connected rank-$r$ graph attaining $\psi(r)$.
Such a graph exists, since cycle counts are integers bounded by $2^r-1$.
Multigraphs have the same extremal value, since uniform subdivision makes them simple without changing rank or cycle count.
Delete leaves and suppress degree-two vertices, preserving rank and cycle count.
For $r\ge2$, the resulting connected multigraph has minimum degree at least three.

Replace each vertex of degree $d\ge4$ by a trivalent tree with $d$ external incidences and $d-2$ vertices.
This adds $d-3$ vertices and $d-3$ edges, so it preserves rank.
Every original cycle has a unique simple lift through the replacement trees.
Contracting those trees recovers the original cycle, so the lift is injective.
The resulting cubic multigraph still attains $\psi(r)$.

Its order is $2r-2$.
It has a cycle $W$ of length $O(\log r)$ by the elementary breadth-first-search girth bound.
A loop or a parallel pair already gives such a cycle.
Otherwise, a girth greater than $2k+1$ would force more than $2^k$ vertices in a depth-$k$ breadth-first tree.

Take $P=V(W)$ and apply Theorem~\ref{thm:two-scale-polynomial-forest} directly.
For any edge $e$ of $W$, map each cycle $C$ containing $e$ to the two even sets
\[
 C,\qquad C\mathbin\triangle W.
\]
All these sets are distinct, since toggling $W$ changes the status of $e$.
Their restrictions to the graph outside $P$ are forests.
Indeed, removing $P$ removes at least one vertex from each cycle $C$, and toggling $W$ changes no outside edge.
Hence
\[
 2c_e\le\Pr(\mathcal F)2^r,
 \qquad c_e\le2^{r-2}+2^{r-1}r^{-\kappa}.
\]
Deleting $e$ preserves connectedness and reduces rank by one.
The same edge recurrence and geometric summation as before give the stated bound.
\end{proof}

\section{Parallel links and an endpoint-scale unrestricted error}
\label{sec:parallel-links-endpoint}

Fable's parallel-link lemma removes the forest excess from the link bound.
This section combines that lemma with two-scale restoration and a separate spectral estimate inside each clean block.
The endpoint parameters below give a relative error $O((\log r)^6/\sqrt r)$.
The argument is mathematical, with independent audits of its constituent estimates.

All logarithms in this section have base two unless written as $\ln$.

\begin{lemma}[Fable's parallel-link estimate]
\label{lem:endpoint-parallel-links}
Let $A,B$ be disjoint connected vertex sets in a connected multigraph $\Gamma$.
Let $P$ be a vertex set, and let $\xi$ be a uniform even subgraph.
Suppose its restriction outside $P$ is a forest with probability greater than one half.
For every integer $t\ge0$, at most $2^{t+1}$ components of $\Gamma-(A\cup B)$ satisfy all four conditions below:
\begin{enumerate}
\item They have an edge to each of $A,B$.
\item They avoid $P$.
\item They contain a cycle.
\item Their boundary has at most $t$ edges.
\end{enumerate}
\end{lemma}

\begin{proof}
If there is no component joining $A$ to $B$, the assertion is immediate.
Otherwise retain one such component $Q_0$.
The set $A\cup B\cup Q_0$ is connected.
For any two other complementary components $Q_i,Q_j$, the graph $\Gamma-Q_i-Q_j$ is connected.
It contains the retained connected set, and every other complementary component has an edge to that set.
No edge joins $Q_i$ and $Q_j$.

Thus the two regions have exactly one link.
The cycle-space pair law makes their zero-boundary events independent.
Conditional on both zero-boundary events, their internal even edge sets are independent and uniform.
This also follows from the direct product of the cycle spaces after deleting the two boundaries.

For each eligible component other than $Q_0$, let $E_i$ mean that its boundary is zero and its internal even edge set is nonzero.
The events $E_i$ are pairwise independent.
The zero-boundary condition imposes at most $t$ linear equations, so its probability is at least $2^{-t}$.
The cyclic internal graph has a nonzero even edge set with conditional probability at least one half.
Consequently $p_i:=\Pr(E_i)\ge2^{-t-1}$.

If there are $m$ such components, put $N=\sum_i\mathbf1_{E_i}$ and $\nu=\mathbb EN$.
Pairwise independence gives
\[
 \mathbb EN^2\le\nu^2+\nu,\qquad
 \Pr(N=0)\le\frac1{1+\nu},\qquad
 \nu\ge m2^{-t-1}.
\]
The forest event implies $N=0$.
Its probability exceeds one half, so $m<2^{t+1}$.
Including $Q_0$, if it was eligible, gives at most $2^{t+1}$ components.
\end{proof}

We recall the walk and moment estimates needed below.
Let $J$ have maximum degree three and girth greater than $D$.
Put $s=\lfloor(D-1)/2\rfloor$, and let $\mathcal C$ contain its cycles of length at most $L$.
For $a_C=2^{-|C|}$ and $\mu=\sum_Ca_C$, the local walk bound gives
\begin{equation}\label{eq:endpoint-local-weight}
 \sum_{C\ni v}a_C\le L2^{-s}.
\end{equation}
The repeated-walk estimate and the even spectral trace estimate give
\begin{equation}\label{eq:endpoint-trace-estimates}
 \mu\le\sum_{\ell\le L}\frac{\operatorname{tr}B_J^\ell}{2\ell2^\ell}
 \le(1+\tfrac32L^22^{-s})\mu,
 \qquad
 2^{-\ell}\operatorname{tr}B_J^\ell
 \ge\left(1-\frac{d_J}{2v_J}\right)^\ell-3v_J2^{-\ell/2}
\end{equation}
for even $\ell$, where $d_J=\sum_v(3-\deg_Jv)$ and $d_J<v_J$.
These are the established non-backtracking walk estimates from Section~\ref{sec:polynomial-foundations}.

Suppose $J=\Gamma-(P\cup S)$, where $S$ comes from a maximal disjoint packing of cycles of length at most $D$ outside $P$.
Write $c=\operatorname{comp}(\Gamma[P])$.
The usual link decomposition, with Lemma~\ref{lem:endpoint-parallel-links} replacing its cyclic-region term, gives
\begin{equation}\label{eq:endpoint-link-estimate}
 \ell(C,C')
 \le c+\frac{2L}{t+1}+2^{t+1}+578(L/D)^2
\end{equation}
whenever $8\le D\le L$ and $2(t+2\lfloor D/8\rfloor)\le D$.
Components meeting $P$ contribute at most $c$.
Components with more than $t$ boundary edges contribute at most $2L/(t+1)$.
The new lemma handles the other cyclic components.
Small tree components and direct edges are bounded by the unchanged girth argument.

\begin{lemma}[The refined second moment]
\label{lem:endpoint-two-link-moment}
Suppose every disjoint pair in $\mathcal C$ has at most $b$ links.
Put
\[
 \eta=\tfrac32\,2^bL^32^{-s}.
\]
Let $W_2$ be the total weight $a_Ca_{C'}$ of ordered disjoint pairs with exactly two links.
If the forest probability exceeds $1/2+\varepsilon$, then
\begin{equation}\label{eq:endpoint-moment-conclusions}
 W_2\ge\mu^2-(1+\eta)\mu,
 \qquad
 \mu\le\frac{1+\eta}{4\varepsilon}.
\end{equation}
The first inequality is strict when $\mu>0$.
\end{lemma}

\begin{proof}
If $\mu=0$, both conclusions hold immediately.
Assume $\mu>0$.
For each $C$, let $I_C$ mean that it is a component of $\xi$.
Put $p_C=\Pr(I_C)$ and $X_C=(a_C/p_C)\mathbf1_{I_C}$.
The cycle-component equations give $p_C\ge a_C$, so
\[
 \mathbb EX_C=a_C,\qquad \mathbb EX_C^2\le a_C.
\]
Distinct intersecting cycles have zero joint contribution.
For disjoint cycles, the pair law gives
\[
 \mathbb E(X_CX_{C'})=a_Ca_{C'}2^{\ell(C,C')-1}.
\]
The exceptional-partner bound and~\eqref{eq:endpoint-local-weight} make the total contribution from pairs with at least three links at most $\eta\mu$.
Thus, for $Z=\sum_CX_C$,
\[
 \mathbb EZ^2\le\mu^2+W_2+(1+\eta)\mu
 \le2\mu^2+(1+\eta)\mu.
\]
The forest event implies $Z=0$.
Cauchy--Schwarz gives
\[
 \mu^2\le\Pr(Z>0)\mathbb EZ^2
 <(1/2-\varepsilon)\mathbb EZ^2.
\]
The first displayed moment bound yields the strict lower bound on $W_2$.
The second yields the upper bound on $\mu$.
\end{proof}

\begin{lemma}[Cycle supply in one clean block]
\label{lem:endpoint-block-supply}
Fix $A=6e^{168}$.
Put $x=\log r$, $L=\lfloor Ax\rfloor$, and $s_0=\lfloor(D_0-1)/2\rfloor$.
Suppose $L^22^{-s_0}=o(1)$.
Let $J$ be a simple graph of maximum degree three, girth greater than $D_0$, and order $v$ satisfying
\[
 8L^2\le v\le3r,
 \qquad d_J\le2L.
\]
For sufficiently large $r$, its cycles of length at most $L$ have total weight at least one.
The threshold is uniform over all such graphs $J$.
\end{lemma}

\begin{proof}
For every even $\ell$ with $6x\le\ell\le L$, Bernoulli's inequality gives
\[
 \left(1-\frac{d_J}{2v}\right)^\ell
 \ge1-\frac{Ld_J}{2v}\ge\frac78.
\]
The nonreal spectral error in~\eqref{eq:endpoint-trace-estimates} is at most $9r^{-2}$.
The even harmonic sum over this interval tends to $\tfrac14\ln(A/6)$.
The repeated-walk factor tends uniformly to one.
Therefore the total cycle weight is at least
\[
 \left(\frac7{32}-o(1)\right)\ln(A/6)>1.
\]
No connectedness or eigenvalue multiplicity estimate is required.
\end{proof}

\begin{theorem}[An endpoint-scale forest error]
\label{thm:parallel-links-endpoint-forest}
Let $\Gamma$ be connected and subcubic, with rank $r$.
Let $P$ contain every vertex of degree less than three.
Put $x=\log r$ and $z=\log x$, and assume
\[
 |P|\le r/x^2,
 \qquad \operatorname{comp}(\Gamma[P])\le x/z^2.
\]
For a uniform even subgraph, let $\mathcal F$ be the event that its restriction to $\Gamma-P$ is a forest.
There is an absolute constant $M$ such that, for all sufficiently large $r$,
\begin{equation}\label{eq:endpoint-forest-log-six}
 \Pr(\mathcal F)\le\frac12+M\frac{x^6}{\sqrt r}.
\end{equation}
Separately, the more conservative parameter choice gives
\begin{equation}\label{eq:endpoint-forest-log-seven}
 \Pr(\mathcal F)\le\frac12+\frac{x^7}{\sqrt r}
\end{equation}
for all sufficiently large $r$.
\end{theorem}

\begin{proof}
If $P$ is empty, the forest event contains only the empty even edge set, so its probability is $2^{-r}$.
Assume $P$ is nonempty and fix one connected component $P_0$ as an anchor.
Keep $A=6e^{168}$, and put
\[
 L=\lfloor Ax\rfloor,
 \quad t=\lfloor z-2\log z\rfloor,
 \quad D_0=\left\lceil\frac{8Ax}{z}\right\rceil.
\]
We prove both bounds with the following two parameter rows:
\[
\begin{array}{c|c|c|c}
 &a&B&\varepsilon\\ \hline
 \text{conservative}&5&0&x^7/\sqrt r\\
 \text{endpoint}&4&\lceil\log(10^6A)\rceil&Mx^6/\sqrt r
\end{array}
\qquad
 D_1=\lfloor x-az-B\rfloor.
\]
The constant $M\ge1$ is chosen at the end, independently of $r$.
Suppose, for a contradiction, that $\Pr(\mathcal F)>1/2+\varepsilon$.
For large $r$, one has $8\le D_0<D_1\le L$ and $t\ge1$.

Choose a maximal disjoint family $\mathcal Q_0$ of cycles of length at most $D_0$ outside $P$.
Extend it to a maximal disjoint family $\mathcal Q_1$ for length $D_1$.
Let $S_j$ be their vertex unions and put $J_j=\Gamma-(P\cup S_j)$.
Then $S_0\subseteq S_1$, and $J_j$ is simple with girth greater than $D_j$.

Each chosen cycle is a cyclic region with boundary at most its length.
The exponential small-cut bound gives
\[
 |\mathcal Q_j|<8D_j2^{D_j}\ln(2/\varepsilon).
\]
Since $M\ge1$, both parameter rows have $\ln(2/\varepsilon)\le x$ for large $r$.
Consequently,
\begin{equation}\label{eq:endpoint-deletion-sizes}
 |S_0|\le8D_0^22^{D_0}x=r^{o(1)},\qquad
 |S_1|\le8\,2^{-B}\frac r{x^{a-3}}.
\end{equation}

Write $v_1=|V(J_1)|$ and $d_1=\sum_{v\in J_1}(3-\deg_{J_1}v)$.
The degree sum in $\Gamma$ gives $|V(\Gamma)|=2r+O(|P|)$, and hence $v_1=2r+o(r)$.
Also $d_1\le3(|P|+|S_1|)$.
For the conservative row, $Ld_1=o(v_1)$.
For the endpoint row, $v_1\ge r$ for large $r$, and
\[
 \frac{Ld_1}{2v_1}
 \le\frac{3A}{2x}+12A2^{-B}<\frac1{100}.
\]
Thus the usual even trace sum over $6x\le\ell\le L$ supplies cycle weight
\begin{equation}\label{eq:endpoint-initial-mass}
 \mu_1:=\sum_{\substack{C\subseteq J_1\\|C|\le L}}2^{-|C|}
 \ge\left(\frac{99}{400}-o(1)\right)\ln(A/6)>20.
\end{equation}
Here $L^22^{-s_1}=o(1)$ controls repeated walks, with $s_j=\lfloor(D_j-1)/2\rfloor$.

At either girth scale, the hypotheses of~\eqref{eq:endpoint-link-estimate} hold for large $r$.
The components meeting $P$ contribute at most $x/z^2$.
Furthermore,
\[
 \frac{2L}{t+1}=(2A+o(1))\frac{x}{z},\quad
 2^{t+1}\le\frac{2x}{z^2},\quad
 578(L/D_j)^2\le578(z/8)^2.
\]
It follows that every link count is at most
\[
 b_r:=3Ax/z.
\]
Since $s_0\ge4Ax/z-2$, the exceptional-pair error satisfies
\[
 \eta_j\le O(L^32^{-Ax/z})=o(1)
 \qquad(j=0,1).
\]
Lemma~\ref{lem:endpoint-two-link-moment} therefore gives
\begin{equation}\label{eq:endpoint-upper-weight}
 \mu_0\le\frac{1+o(1)}{4\varepsilon},
 \qquad W_{2,1}>\mu_1^2-(1+o(1))\mu_1.
\end{equation}

We next obtain a nested chain at the larger girth scale.
For disjoint cycles $C,C'$, call $C'$ far behind $C$ if it lies outside the $P_0$-component of $\Gamma-V(C)$.
This relation is antisymmetric.
A path from $P_0$ to its first vertex on $C$ avoids that far component and hence avoids $C'$.
Connectedness then puts all of $C$ in the $P_0$-component after deleting $C'$.
Thus ordered far pairs have total weight at most $\mu_1^2/2$.
The weight of pairs joined by an edge is at most $\mu_1L^22^{-s_1}=o(\mu_1)$.

Equations~\eqref{eq:endpoint-initial-mass} and~\eqref{eq:endpoint-upper-weight} therefore give a cycle $C$ with partners of total weight at least eight satisfying three conditions.
They lie in the $P_0$-component of $H:=\Gamma-V(C)$, have exactly two links to $C$, and have no edge to $C$.
For each partner, one of the two exterior linking components avoids $P_0$.
It contains a neighbor $u$ of $C$, and the partner separates $u$ from $P_0$ in $H$.
There are at most $L$ choices of $u$.
For one choice, the separating cycles have total weight at least $8/L$.

Take a maximal vertex-disjoint subfamily of these cycles, of size $T$.
Its vertices meet every separating cycle in the family.
The local weight estimate gives
\begin{equation}\label{eq:endpoint-chain-size}
 T\ge\frac{8\,2^{s_1}}{L^3}
 \ge\frac{2^{1-B/2}}{A^3}\frac{\sqrt r}{x^{a/2+3}}.
\end{equation}
The cycles are nested between $u$ and $P_0$ in $H$.
Indeed, two disjoint connected separators cannot cross one another's sides.
Rerouting a $u$--$P_0$ path through the first separator shows that the second lies entirely on one side.
The separation property then gives a consistent order.
Write the ordered cycles as $C_1,\ldots,C_T$, and let $U_i$ be the $u$-component of $H-V(C_i)$.
Then $U_i\subseteq U_{i+1}$ and $C_i\subseteq U_{i+1}$.

Put $q=\lceil16A^2x\rceil+2$.
Choose disjoint intervals $[a',b']$ of indices with $b'-a'=q$, leaving a separator between successive intervals.
For each interval, use the full vertex set
\[
 V(a',b')=U_{b'}\setminus(U_{a'}\cup V(C_{a'})).
\]
It contains all $q-1$ intermediate cycles.
Thus, since $D_1+1\ge x/2$ for large $r$,
\[
 |V(a',b')|\ge(q-1)(D_1+1)\ge8L^2.
\]
Every edge of $H$ leaving this set ends on $C_{a'}$ or $C_{b'}$.
Its boundary in $H$ is at most $2L$.
The chosen full sets are pairwise disjoint.

Discard blocks meeting a neighbor of $C$, a vertex of $S_0$, or a vertex of $P$.
At most $L+|S_0|+c$ blocks are discarded.
For the last term, each connected component of $\Gamma[P]$ can meet at most one block.
It cannot cross a block boundary because every such boundary ends on a separator cycle disjoint from $P$.

Each retained block avoids $P\cup S_0$ and has no neighbor in either set.
An edge to either set would leave the block and end on a separator cycle, which avoids both sets.
The discarded neighbors of $C$ remove the only edges missing from this boundary description in $\Gamma$.
Thus the induced graph on each retained block has girth greater than $D_0$, and degree deficit at most $2L$.
Its order is at most $|V(\Gamma)|\le3r$.
Lemma~\ref{lem:endpoint-block-supply} gives cycle weight at least one in each retained block.

The losses in~\eqref{eq:endpoint-deletion-sizes} are negligible compared with the number of blocks.
Using~\eqref{eq:endpoint-chain-size} and $q+1\le18A^2x$, for large $r$ their number $m$ satisfies
\begin{equation}\label{eq:endpoint-clean-block-count}
 m\ge c_B\frac{\sqrt r}{x^{a/2+4}},
 \qquad c_B:=\frac{2^{-B/2}}{64A^5}.
\end{equation}
Cycles in different blocks are distinct cycles of $J_0$.
Consequently $\mu_0\ge m$.

For the conservative row, this contradicts
\[
 c_0\frac{\sqrt r}{x^{13/2}}
 \le\mu_0\le(1+o(1))\frac{\sqrt r}{4x^7}.
\]
For the endpoint row, choose the absolute constant
\[
 M=\frac1{c_B}=64A^5 2^{B/2}.
\]
Then it contradicts
\[
 c_B\frac{\sqrt r}{x^6}
 \le\mu_0\le(1+o(1))\frac{c_B}{4}\frac{\sqrt r}{x^6}.
\]
This proves both forest estimates.
\end{proof}

\begin{corollary}[Unrestricted endpoint-scale cycle error]
\label{cor:parallel-links-endpoint-cycles}
There is an absolute constant $M$ such that
\[
 \psi(r)\le2^{r-1}+(M+o(1))\frac{2^r(\log r)^6}{\sqrt r}.
\]
In particular, for every fixed $\kappa<1/2$,
\[
 \psi(r)\le2^{r-1}+(1+o(1))2^r r^{-\kappa}.
\]
\end{corollary}

\begin{proof}
Use the direct cubic-extremizer transfer of Proposition~\ref{prop:two-scale-direct-cycle-transfer}.
A shortest cycle gives a connected marked set of size $O(\log r)$, satisfying the new hypotheses.
The cycle-flip injection bounds some cyclic edge count by
\[
 c_e\le2^{r-2}+\frac M2\frac{2^r(\log r)^6}{\sqrt r}.
\]
Deleting that edge gives the same rank recurrence as before.
The geometric tail satisfies
\[
 \sum_{j\le r}2^{j-1}\frac{(\log j)^6}{\sqrt j}
 =(1+o(1))\frac{2^r(\log r)^6}{\sqrt r}.
\]
This proves the first bound.
For each fixed $\kappa<1/2$, one has $(\log r)^6/\sqrt r=o(r^{-\kappa})$.
The second bound follows.
\end{proof}

\section{A sharper packing estimate reduces the logarithmic factor}
\label{sec:parallel-links-log-five-half}

The stronger branch of the exponential cut proof improves the endpoint parameters by one half-power of $\log r$.
We record this refinement separately from the preceding audited parameter choices.

\begin{lemma}[Packing bound above one half]
\label{lem:endpoint-sharp-packing}
Suppose $k$ disjoint connected cyclic regions avoid the marked set and have boundaries at most an integer $d\ge2$.
If the forest probability exceeds $1/2+\varepsilon$, with $0<\varepsilon<1/2$, then
\[
 k<2^{d+1}\max\left\{5d+1,\ln\frac1{2\varepsilon}\right\}.
\]
\end{lemma}

\begin{proof}
The proof of Theorem~\ref{thm:sharper-exponential-cuts} gives the stronger implication
\[
 k\ge2(5d+1)2^d
 \quad\Longrightarrow\quad
 \Pr(\mathcal F)-\frac12
 \le\frac12\exp\left(-\frac{k}{2^{d+1}}\right).
\]
The cyclic-region version follows from the same quotient and conditional internal-cycle-space argument.
If the threshold fails, the first term in the maximum gives the asserted bound.
Otherwise, compare the displayed exponential estimate with $\varepsilon$ and take logarithms.
\end{proof}

\begin{theorem}[The improved endpoint forest bound]
\label{thm:parallel-links-log-five-half-forest}
Under the hypotheses of Theorem~\ref{thm:parallel-links-endpoint-forest}, there is an absolute constant $M$ such that
\[
 \Pr(\mathcal F)
 \le\frac12+M\frac{(\log r)^{11/2}}{\sqrt r}
\]
for all sufficiently large $r$.
One may take
\[
 A=6e^{168},\qquad B=\lceil\log(10^6A)\rceil,
 \qquad M=64A^5 2^{B/2}.
\]
\end{theorem}

\begin{proof}
Keep $x=\log r$, $z=\log x$, and the parameters
\[
 L=\lfloor Ax\rfloor,\quad
 t=\lfloor z-2\log z\rfloor,\quad
 D_0=\left\lceil\frac{8Ax}{z}\right\rceil.
\]
Replace only the large girth scale and the assumed excess by
\[
 D_1=\lfloor x-3z-B\rfloor,
 \qquad \varepsilon=Mx^{11/2}/\sqrt r.
\]
Construct the two nested disjoint short-cycle packings as before.
For both scales, $D_j\le x$ and $\ln(1/(2\varepsilon))\le x$ for large $r$.
Lemma~\ref{lem:endpoint-sharp-packing} therefore gives
\[
 |\mathcal Q_j|<12x2^{D_j},\qquad
 |S_j|<12x^22^{D_j}.
\]
In particular,
\[
 |S_0|=r^{o(1)},\qquad
 |S_1|\le12\,2^{-B}r/x.
\]

As before, $v_1=2r+o(r)$ and $d_1\le3(|P|+|S_1|)$.
For large $r$, the fixed spectral loss satisfies
\[
 \frac{Ld_1}{2v_1}
 \le\frac{3A}{2x}+18A2^{-B}<\frac1{100}.
\]
The initial cycle supply therefore remains $\mu_1>20$.
The link bound is unchanged, since it contains no forest-excess term.
It gives at most $3Ax/z$ links at either scale.
Thus the exceptional-pair error remains $O(L^32^{-Ax/z})=o(1)$.
Both conclusions of Lemma~\ref{lem:endpoint-two-link-moment} still apply.

The large-scale girth now gives
\[
 2^{s_1}\ge2^{-B/2-2}\frac{\sqrt r}{x^{3/2}}.
\]
The same weighted averaging and maximal-disjointness argument produce a nested chain of size at least
\[
 T\ge\frac{2^{1-B/2}}{A^3}\frac{\sqrt r}{x^{9/2}}.
\]
Use separator intervals with $q=\lceil16A^2x\rceil+2$ gaps and the same full block definition.
Each block has at least $8L^2$ vertices and boundary at most $2L$.
Discarding blocks meeting $P$, $S_0$, or a neighbor of the anchor cycle loses only $r^{o(1)}$ blocks.
The retained block count is at least
\[
 m\ge c_B\frac{\sqrt r}{x^{11/2}},
 \qquad c_B=\frac{2^{-B/2}}{64A^5}.
\]
The one-block spectral supply gives $\mu_0\ge m$.
The inverse moment bound instead gives
\[
 \mu_0\le\frac{1+o(1)}{4\varepsilon}
 =(1+o(1))\frac{c_B}{4}\frac{\sqrt r}{x^{11/2}},
\]
since $M=1/c_B$.
These inequalities contradict each other for large $r$.
\end{proof}

\begin{corollary}[Unrestricted cycle bound]
\label{cor:parallel-links-log-five-half-cycles}
There is an absolute constant $M$ such that
\[
 \psi(r)
 \le2^{r-1}+(M+o(1))\frac{2^r(\log r)^{11/2}}{\sqrt r}.
\]
\end{corollary}

\begin{proof}
Apply the direct cubic-extremizer transfer and its geometric rank summation, now with exponent $11/2$ on the logarithm.
\end{proof}


\section{Verification and remaining questions}
The proof is independent of the terminal-path saturation theorem and the separate results for cuts of sizes two, three, and four.
Its inputs are the exponential cut estimate, the parallel-link lemma, and the walk and cycle-space identities proved here.

Exact finite checks independently test the cut-space decompositions and block conditional laws.
A separate exact checker tests 9,759 disjoint cycle pairs, 200 weighted antisymmetry inequalities, and 521 clean restored blocks in 160 chains.
Every retained block has a strict deficit reduction after restoration.
Independent numerical spectral checks test the operator identities and eigenvalue counts.
The final endpoint proof uses a separate spectral supply estimate in each block, so it does not need the earlier multiplicity argument.
The parallel-link identities also have an independent exact check on 2,228 component families.
These checks support individual identities; they do not replace the asymptotic proof.

The constants and the associated rank threshold are very large.
No useful finite-rank numerical estimate is claimed from them.
The proposed bound $\psi(r)-2^{r-1}=O(r2^{r/2})$ remains open in this research programme.
The original manuscripts and existing root formalization remain unchanged.
The new supporting Lean modules are separate, and the full new theorem still requires formal assembly.
\end{document}
